% Paratechnical Data Science Handbook
% Chapter: Using AI and LLM Tools as Part of Technical Practice
% Voice: textbook / O’Reilly style (concise, structured, verification-forward)
% Target length: ~4000 words
% Assumes you have defined the verbatim environment in the preamble.

\chapter{Using AI Tools}
\label{ch:ai-llm}

\section*{Purpose}
AI assistants can reduce friction in programming and data science. They can draft documentation, summarize error messages, propose alternative implementations, and generate boilerplate. They can also produce plausible but incorrect guidance. This chapter treats AI assistance as one input to a disciplined workflow. The goal is not to get answers faster. The goal is to produce correct, reproducible, and safe work with less unnecessary effort.

The governing rule is simple:

\begin{quote}
\textbf{AI can propose. You must verify.}
\end{quote}

Verification means using primary documentation, running small experiments, and adding tests. Those practices remain your responsibility, even if an assistant produced the first draft.

\section*{Learning objectives}
By the end of this chapter, you should be able to:
\begin{enumerate}[noitemsep]
  \item Identify high-value AI uses and common failure modes.
  \item Apply a risk-based policy to decide how much verification is required.
  \item Write prompts that produce usable how-to guidance, not vague advice.
  \item Validate AI outputs using primary sources, minimal experiments, and automated tests.
  \item Avoid privacy and security mistakes when sharing context.
  \item Use AI in coursework and collaboration without undermining learning or integrity.
\end{enumerate}

\section{What an AI assistant is doing}
Most AI assistants in technical settings are built on large language models (LLMs). An LLM generates text that fits patterns in its training data and in your prompt. That makes it useful for drafting and transformation tasks, such as turning rough notes into a README or converting a stack trace into a troubleshooting plan. It does not guarantee factual correctness.

A practical interpretation is to treat AI output as a draft produced by a fast collaborator who can be wrong. Your job is to turn drafts into reliable artifacts.

\subsection{What AI is usually good at}
LLM assistance is strongest when the task is primarily about language, structure, or recall of common patterns:
\begin{itemize}[noitemsep]
  \item Drafting outlines, READMEs, docstrings, issue templates, and lab instructions.
  \item Rewriting an explanation for a novice reader.
  \item Producing boilerplate code patterns (argument parsing, logging setup) that you then review.
  \item Suggesting search terms, likely documentation sections, and common failure modes.
  \item Generating checklists that you then shorten and validate.
\end{itemize}

\subsection{What AI is not reliable at}
These failures appear frequently enough that you should plan around them:
\begin{itemize}[noitemsep]
  \item \textbf{Hallucinated details}: invented flags, APIs, or citations.
  \item \textbf{Version mismatch}: correct advice for a different version of a tool than yours.
  \item \textbf{Hidden prerequisites}: missing steps such as activating an environment, changing directories, or installing a system dependency.
  \item \textbf{Silent logic bugs}: code that runs but is wrong on edge cases.
  \item \textbf{Unsafe defaults}: suggestions that weaken security or increase destructive behavior.
\end{itemize}

You should assume any of these errors are possible. This assumption leads to good habits: request sources, run small tests, and prefer minimal changes.

\section{A risk-based verification policy}
Not all tasks have the same consequences if something goes wrong. Use risk level to decide how careful you need to be.

\subsection{Low risk: drafting and formatting}
Low-risk tasks are those where an error is easy to spot and easy to correct:
\begin{itemize}[noitemsep]
  \item Rewrite a paragraph for clarity.
  \item Draft a README skeleton.
  \item Produce a checklist or template.
\end{itemize}

Verification: read for accuracy, completeness, and alignment with your assignment or project requirements.

\subsection{Medium risk: technical guidance you can test quickly}
Medium-risk tasks involve technical claims but can be validated with small experiments:
\begin{itemize}[noitemsep]
  \item Interpreting a stack trace and proposing checks.
  \item Suggesting commands to inspect an environment.
  \item Drafting a unit test scaffold.
  \item Proposing a refactor that you can validate with tests.
\end{itemize}

Verification: run the checks, confirm in primary documentation, and add tests or assertions.

\subsection{High risk: destructive commands, security changes, and sensitive data}
High-risk tasks require strict guardrails:
\begin{itemize}[noitemsep]
  \item Commands that delete, overwrite, or move files recursively.
  \item Privilege escalation (\texttt{sudo}) and broad permission changes (\texttt{chmod 777}).
  \item Network and authentication configuration (SSH keys, VPNs, firewall rules).
  \item Any handling of protected, confidential, or regulated data.
\end{itemize}

Verification: consult primary documentation; prefer supervised help (instructor/TA); test in a sandbox; do not execute commands you do not understand.

\subsection{A non-negotiable warning for destructive commands}
Treat these as stop signs until you have confirmed paths, backups, and intent:

\begin{verbatim}
# Examples of destructive patterns (do not run blindly)
# rm -rf <path>              # macOS/Linux: recursive delete
# del /s <path>              # Windows: recursive delete
# mv <source> <dest>         # can overwrite depending on flags/context
# chmod -R 777 <path>        # broad permissions (usually wrong)
# sudo <anything>            # elevated privileges
\end{verbatim}

If an assistant suggests one of these as a fix, pause. Ask: What exactly will this change? What evidence says this is the right change? How do I undo it?

\section{The assistive loop: a workflow that forces evidence}
A reliable workflow is an assistive loop: you use AI to propose an approach, then you validate it with evidence.

\subsection{Step 1: write a short specification}
Include:
\begin{itemize}[noitemsep]
  \item Goal (one sentence).
  \item Context (OS, Python version, environment name, relevant package versions).
  \item Inputs (code snippet, command, file path, sample data).
  \item Observed behavior (exact error/output).
  \item Expected behavior.
  \item What you tried and what happened.
\end{itemize}

This structure is also the structure of a good technical question.

\subsection{Step 2: request alternatives and checkpoints}
Avoid prompts that ask for a single answer. Ask for:
\begin{itemize}[noitemsep]
  \item 2 to 4 plausible root causes.
  \item A short decision tree of checks.
  \item A step-by-step plan with verification checkpoints.
\end{itemize}

This reduces the chance you follow one incorrect path.

\subsection{Step 3: validate claims in primary documentation}
If the output depends on external truth (API behavior, command flags), verify it in official documentation for your version. If the assistant cannot point to a primary source, treat the claim as tentative.

\subsection{Step 4: run a minimal experiment}
Convert advice into the smallest experiment that can succeed or fail. Change one factor at a time. Record results.

\subsection{Step 5: lock in the outcome}
After a fix, add one of:
\begin{itemize}[noitemsep]
  \item a unit test,
  \item an assertion (shape, type, range),
  \item a checklist item in project documentation.
\end{itemize}

This is how you prevent regression.

\section{Prompt patterns that produce usable how-to guidance}
Prompts work best when they specify output format and constraints. The goal is to be unambiguous.

\subsection{Pattern A: decision tree diagnosis}
\begin{fullcode}
# "I am seeing: <symptom or exact error>. Provide a decision tree with 8 checks.
# Each check must be a concrete command or observation, and you must say what
# each outcome implies. Keep it specific to <OS> and <tool/version>."
\end{fullcode}

\subsection{Pattern B: minimal reproducible example reduction}
\begin{fullcode}
# "Here is my current minimal failing snippet. Reduce it further if possible.
# Replace real data with synthetic data. Output:
# (1) reduced code, (2) what you removed and why, (3) how to verify failure."
\end{fullcode}

\subsection{Pattern C: test-first scaffold}
\begin{fullcode}
# "Write 4 pytest tests for the intended behavior below (include 2 edge cases).
# Then propose an implementation that satisfies the tests. State assumptions."
\end{fullcode}

\subsection{Pattern D: documentation rewrite without inventing steps}
\begin{fullcode}
# "Rewrite these notes into a how-to guide with sections: Purpose, Prerequisites,
# Steps, Verify, Troubleshooting. Do not invent commands I did not provide.
# If prerequisites are missing, list questions under 'Prerequisites'."
\end{fullcode}

\subsection{Pattern E: safe command review}
Use when you want to understand a command before running it.
\begin{fullcode}
# "Explain what this command does, what it changes on disk, and how to undo it:
# <command>
# Then give a safer alternative or a dry run option if available."
\end{fullcode}

\subsection{Pattern F: request a patch instead of a rewrite}
For code changes, request a minimal diff rather than a full replacement.
\begin{fullcode}
# "Here is the current function. Provide a minimal patch (diff-style) to fix
# <specific bug>. Do not refactor unrelated code. Explain how to test the fix."
\end{fullcode}

\section{Using AI to improve technical questions}
Good questions produce good answers. AI can help you edit and structure questions before you ask a human.

\subsection{Rewrite into a standard template}
Ask the assistant to convert a messy description into:
\begin{itemize}[noitemsep]
  \item Goal
  \item Expected
  \item Actual
  \item Reproduction steps
  \item Context (versions, OS)
  \item What you tried
\end{itemize}

Then audit the result for correctness. Do not let an assistant improve a question by inventing details.

\subsection{Generate a context checklist}
If you do not know what context matters, ask for a checklist. Common context fields include:
\begin{itemize}[noitemsep]
  \item OS and version
  \item Python version and interpreter path
  \item Environment manager (conda, venv) and active environment name
  \item Package versions (\texttt{pip show} or \texttt{conda list})
  \item Working directory and relevant file paths
\end{itemize}

Report unknowns as unknowns rather than guessing.

\section{Using AI in debugging: hypotheses, checks, and minimal diffs}
Use AI to propose hypotheses and checks. Keep control of the debugging loop.

\subsection{Import errors: identify the interpreter and environment}
Most import problems are environment problems. Confirm which Python is running and where packages are installed.

\begin{fullcode}
# Interpreter and environment checks
python --version
python -c "import sys; print(sys.executable)"

# macOS/Linux:
which python

# Windows:
where python

# Package presence:
pip show <package>
conda list <package>
\end{fullcode}

Interpretation:
\begin{itemize}[noitemsep]
  \item If \texttt{sys.executable} is not in the environment you expect, you are using the wrong interpreter.
  \item If \texttt{pip show} reports a package under a different Python than \texttt{sys.executable}, you installed into the wrong environment.
\end{itemize}

Ask the assistant to propose a diagnosis tree, but do not accept environment-changing commands without understanding them.

\subsection{Data ingestion bugs: inspect the intermediate}
For many data issues, the right move is to inspect intermediate representations.

Example: a CSV loads into one column instead of many. Plausible causes include the wrong delimiter, quoting, or encoding artifacts. A useful debugging sequence is:
\begin{enumerate}[noitemsep]
  \item print the first few lines of the file,
  \item test delimiters explicitly,
  \item check column names and types after load.
\end{enumerate}

\begin{fullcode}
# Quick inspection approach (Python)
import pandas as pd

path = "data/input.csv"
with open(path, "r", encoding="utf-8", errors="replace") as f:
    for _ in range(5):
        print(f.readline().rstrip("\n"))

df = pd.read_csv(path)              # default delimiter is comma
print(df.shape)
print(df.columns.tolist())

df_tab = pd.read_csv(path, sep="\t")
print(df_tab.shape)
\end{fullcode}

Use AI to suggest likely causes and checks, but interpret the output yourself. A model cannot see your file contents unless you provide them, and you should avoid pasting sensitive data.

\subsection{Path and working directory confusion}
Many errors are simple path mistakes. Validate where you are before changing code.
\begin{fullcode}
# Working directory checks
# macOS/Linux:
pwd
ls -la

# Windows:
cd
dir
\end{fullcode}

If AI suggests changing or deleting files to solve a path problem, stop. Verify the path and fix the invocation first.

\subsection{Prefer minimal diffs over large rewrites}
A common anti-pattern is accepting a large rewrite because it makes an error disappear. Large changes reduce your ability to identify cause and increase the chance of introducing new bugs. If an assistant proposes a rewrite:
\begin{itemize}[noitemsep]
  \item ask for the smallest change that preserves intent,
  \item apply one change at a time,
  \item verify with a test or a checkpoint.
\end{itemize}

\section{Using AI for documentation and project hygiene}
Documentation is a strong target for AI assistance because outputs are reviewable and verifiable.

\subsection{Draft a README that can be executed}
A README is good only if a reader can follow it. Use AI to draft structure, then validate it by running it from scratch.

Minimum sections:
\begin{itemize}[noitemsep]
  \item Purpose
  \item Setup (environment creation)
  \item How to run
  \item Verify (what success looks like)
  \item Troubleshooting (common failures)
\end{itemize}

Provide the actual commands you use. Ask the assistant to format them. Then test them in a fresh environment. This is often where missing steps are discovered.

\subsection{Troubleshooting sections: symptom, cause, check, fix}
Ask the assistant to propose troubleshooting entries in a strict format:
\begin{itemize}[noitemsep]
  \item Symptom
  \item Likely cause
  \item Check (a command or observable)
  \item Fix (a minimal action)
\end{itemize}

Then validate each fix on your system. Remove entries you cannot verify.

\subsection{Runbooks for repetitive tasks}
For recurring tasks (refreshing a dataset, rebuilding outputs), create a runbook. Include:
\begin{itemize}[noitemsep]
  \item prerequisites and inputs,
  \item step-by-step commands,
  \item verification checks,
  \item rollback or recovery steps.
\end{itemize}

AI can draft a runbook; you must execute it to confirm it works.

\section{Using AI for code: constraints, review, and tests}
AI-generated code is useful when you control scope and require verification.

\subsection{Request small, testable units}
Instead of asking for an entire pipeline, request one function with a clear contract:
\begin{itemize}[noitemsep]
  \item Inputs: types and constraints
  \item Output: type and invariants
  \item Examples: at least one realistic case
  \item Tests: include edge cases
\end{itemize}

This makes review manageable and encourages correct design.

\subsection{Verification ladder}
Use this ladder for any AI-generated code you plan to keep:
\begin{enumerate}[noitemsep]
  \item Read it line by line. You should be able to explain it.
  \item Run a smoke test on tiny inputs.
  \item Add assertions for invariants (shape, type, ranges).
  \item Add unit tests (including edge cases).
  \item Compare outputs to a baseline if replacing existing code.
\end{enumerate}

If you cannot explain the code, do not ship it. ``It works'' is not a sufficient reason to keep code you do not understand, especially in an educational context.

\subsection{Ask for tests before implementations}
A reliable pattern is to ask for tests first. Tests clarify the contract and catch silent errors.

For example, if you are cleaning a column:
\begin{itemize}[noitemsep]
  \item define how empty strings should be treated,
  \item define how whitespace is handled,
  \item define whether case is preserved,
  \item define how non-string values should behave.
\end{itemize}

Then request tests that lock that behavior in.

\subsection{Warnings for file operations in code}
Code that writes files can destroy work. When an assistant proposes code that deletes or overwrites:
\begin{itemize}[noitemsep]
  \item request a dry run mode,
  \item request explicit confirmation before destructive operations,
  \item request that outputs go to a dedicated \texttt{outputs/} directory,
  \item confirm the code uses explicit paths, not implicit working-directory assumptions.
\end{itemize}

\section{Security and privacy hygiene}
Security mistakes often begin as convenience: pasting too much context.

\subsection{Never paste secrets}
Do not paste:
\begin{itemize}[noitemsep]
  \item passwords,
  \item API keys or tokens,
  \item private SSH keys,
  \item institutional credentials,
  \item confidential datasets.
\end{itemize}

If you need help, describe the structure and the error. If you need an example, use synthetic values.

\subsection{Prefer synthetic or summarized data}
If you need help with data issues, share:
\begin{itemize}[noitemsep]
  \item schema (columns and types),
  \item 3 to 5 synthetic rows that match structure,
  \item aggregate statistics (counts, missingness, ranges).
\end{itemize}

Avoid copying raw records when they include personal or sensitive information. If policy or law applies, assume raw disclosure is not allowed.

\subsection{Treat privilege escalation as high risk}
Commands that request elevated privileges require extra caution. If an assistant suggests \texttt{sudo} or broad permission changes:
\begin{itemize}[noitemsep]
  \item consult primary documentation,
  \item confirm the scope of the change,
  \item prefer least-privilege alternatives,
  \item ask a human when you are unsure.
\end{itemize}

\section{Academic integrity and collaboration}
Courses and teams differ in their policies. Follow the policy that governs your work.

\subsection{Disclose AI assistance when required}
When disclosure is expected, keep it concrete:
\begin{itemize}[noitemsep]
  \item ``Used AI assistance to draft README structure and suggest unit test scaffolding; verified commands by running locally; edited output for accuracy.''
\end{itemize}

Avoid vague statements that do not clarify what was assisted and what was verified.

\subsection{Use AI to support learning}
If you receive an answer, convert it into learning:
\begin{itemize}[noitemsep]
  \item Ask for a smaller example.
  \item Ask for assumptions and edge cases.
  \item Predict what will happen if an input changes.
  \item Test those predictions.
\end{itemize}

A useful personal rule is: if you cannot explain it, you do not own it.

\subsection{Collaboration: keep humans in the loop}
In team settings, AI can accelerate drafting, but humans should review:
\begin{itemize}[noitemsep]
  \item Pull requests should explain intent and include tests.
  \item Code review comments should reference evidence (tests, logs, docs).
  \item Generated code should be treated like any other contribution: reviewed, tested, and justified.
\end{itemize}

\section{When AI advice conflicts with your observations}
Conflicts are common. Resolve them with evidence.

\subsection{Resolution protocol}
\begin{enumerate}[noitemsep]
  \item Trust local evidence: what happened when you ran the code.
  \item Confirm versions and environment.
  \item Consult primary documentation for that version.
  \item If still unclear, ask a human with a structured question.
\end{enumerate}

If the assistant cannot provide a verifiable primary source, treat its claim as tentative.

\section{Case studies}
These case studies illustrate how AI fits into a disciplined workflow.

\subsection{Case study 1: notebook kernel mismatch}
Symptom: a package imports in the terminal but fails in Jupyter.

Workflow:
\begin{enumerate}[noitemsep]
  \item Use AI to propose checks: inspect \texttt{sys.executable} in the notebook.
  \item Run the check and confirm the notebook kernel points to the wrong environment.
  \item Switch kernels or install into the correct environment.
  \item Add a note to your README: ``Kernel must be \texttt{<env-name>}.''
\end{enumerate}

AI helped propose checks. Local evidence determined the fix.

\subsection{Case study 2: CSV delimiter confusion}
Symptom: a CSV loads into one column.

Workflow:
\begin{enumerate}[noitemsep]
  \item Ask AI for plausible causes: delimiter, quoting, encoding.
  \item Inspect the first lines of the file; test \texttt{sep} explicitly.
  \item Confirm the correct delimiter and update load code.
  \item Add an assertion: expected columns present.
\end{enumerate}

The fix is reliable because it is backed by inspection and an assertion.

\subsection{Case study 3: Git merge conflict in a notebook}
Symptom: merge conflict in \texttt{analysis.ipynb}.

Workflow:
\begin{enumerate}[noitemsep]
  \item Ask AI to outline options: resolve in an editor, use a notebook diff tool, or rerun to regenerate.
  \item Choose an option consistent with your course and team workflow.
  \item After resolving, add a practice note: avoid committing noisy notebook outputs if your policy allows (see Chapter~\ref{ch:git-github} for Git and version control practices).
\end{enumerate}

AI can propose approaches, but your team's policy determines the acceptable option.

\subsection{Case study 4: unsafe suggested fix}
Symptom: permission error writing to a directory.

The assistant suggests \texttt{sudo} or a broad permission change.

Correct response:
\begin{itemize}[noitemsep]
  \item Stop and identify the correct output directory.
  \item Prefer writing to a user-owned folder (project \texttt{outputs/} under your home directory).
  \item Consult documentation or a TA if you think system permissions are required.
\end{itemize}

High-risk suggestions require human judgment and primary sources.

\section{Exercises}
\begin{enumerate}[noitemsep]
  \item Take a recent error. Ask an assistant for three hypotheses and two checks per hypothesis. Run the checks and record which hypotheses you eliminated.
  \item Draft a structured help request (Goal/Expected/Actual/Repro/Context/What I tried). Ask the assistant to improve clarity without changing facts. Audit the result and then post it.
  \item Ask for a README skeleton for your project. Verify each command in a fresh environment. Add a Verify section and one troubleshooting entry.
  \item Ask for a small function plus unit tests. Run the tests. If any fail, revise until they pass.
  \item Create a personal ``do not paste'' list (credentials, tokens, private data). Keep it near your workstation.
\end{enumerate}

\section{One-page checklist}
\begin{itemize}[noitemsep]
  \item I treat AI output as a proposal and require evidence.
  \item I scale verification to risk.
  \item I request alternatives and checkpoints, not a single answer.
  \item I confirm commands and API behavior in primary documentation.
  \item I run minimal experiments and change one thing at a time.
  \item I add tests or assertions after fixes.
  \item I do not paste secrets or sensitive data.
  \item I disclose AI assistance when required and keep it specific.
\end{itemize}

% \section*{Further reading}